\documentclass[10pt]{article}
\usepackage[margin=0.75in]{geometry}
\usepackage{amsmath,microtype}
\setlength{\emergencystretch}{3em}
\title{Guide: The Correction-Coverage Gap Is Closed}
\author{Collatz proof project}
\date{September 5, 2026}
\begin{document}
\maketitle
\section*{What Lean now proves}
Suppose a positive natural Collatz orbit is unbounded. Then the sum of the
reciprocals of all its successive values converges. This is a theorem about
every hypothetical unbounded orbit, not about almost every starting seed.

The existing reciprocal theorem then gives bounded ideal correction.
The previously unresolved case of an unbounded orbit with divergent
reciprocal sum is now ruled out.

\section*{The key idea}
An unbounded orbit never repeats a value. Its distinct values cannot merge
after the same number of steps, since that would create a repetition.
Most small starting values map into a smaller range after suitably many
steps. Injectivity limits how many of them can belong to one unbounded orbit.

Lean proves that a finite set of $S$ orbit values below $32^m$ satisfies
$8^mS\le2\cdot216^m+243^m$. Splitting the positive integers into shells
from $32^m$ to $32^{m+1}$ turns this into a convergent reciprocal majorant.

\section*{What this adds to the proof program}
Every hypothetical unbounded positive orbit now belongs to the existing
bounded-correction framework. Its inverse drift $2^t/3^{j_t}$ tends to zero,
and its shifted correction divided by its current value tends to zero.
Neither statement supplies a fixed gap in average odd-step density.

The quantitative refinement gives an explicit budget $B_K\to0$: an
unbounded orbit that always stays above $32^K$ has reciprocal sum at most
$B_K$ and correction at most $N(e^{B_K}-1)$. The technical paper gives
the two geometric terms in $B_K$ and their Lean proofs.

This bound does not make the absolute correction small. Lean proves that
the upper envelope, even at $N=32^K$, is at least $(1728/5)27^K$.
That is a limitation of this estimate, not a lower bound on actual errors.
A separate verified barrier puts actual shifted corrections above one.
We therefore still need an arithmetic argument beyond rounding the ideal.

\section*{Why this is not Collatz yet}
Bounded correction is a restriction on an unbounded orbit; it is not an
exclusion of one. Nontrivial cycles also remain a separate unresolved case.
Finite cycles repeat values, so their infinite reciprocal series diverges;
they do not satisfy the new summability hypothesis.

The mathematical implication is related to a quantitative 1999 result of
Garcia and Tal. We make no novelty claim. The main project's proof uses its
own existing bounds and adds no external mathematical axiom.

\section*{Verification}
The Lean build and 120-declaration axiom audit pass, as do three exact finite
controls. Run \texttt{lake build Collatz.OrbitSummability} from \texttt{lean/}.
The technical paper provides the argument and precise remaining gap.
\end{document}
